\documentclass[11pt,a4paper]{article}
\usepackage[utf8]{inputenc}
\usepackage[T1]{fontenc}
\usepackage[french]{babel}
\usepackage{amsmath, amssymb, amsthm}
\usepackage{geometry}
\geometry{margin=2.5cm}

\title{Corrigé - TD 2 : Espaces probabilisés}
\author{MT 331 - INP Esisar / UGA}
\date{Année 2025-2026}

\begin{document}

\maketitle

\section*{Exercice 1}
On note les événements suivants : 
$M_a$ : "la personne est prof de maths", 
$M_e$ : "la personne est médecin", 
[span_0](start_span)$N$ : "la personne est normale"[span_0](end_span). 
[span_1](start_span)On note également $L$ : "le mot est lisible" et $\bar{L}$ : "le mot est illisible"[span_1](end_span).
[span_2](start_span)D'après l'énoncé, on a : $P(M_a) = 0.15$, $P(M_e) = 0.08$ et $P(N) = 0.77$[span_2](end_span). 
[span_3](start_span)Les probabilités conditionnelles d'écrire de façon illisible sont : $P(\bar{L}|M_a) = 0.5$, $P(\bar{L}|M_e) = 0.7$ et $P(\bar{L}|N) = 0.3$[span_3](end_span).

\subsection*{1. Probabilité que le mot soit illisible}
[span_4](start_span)\textbf{Raisonnement :} La famille $\{M_a, M_e, N\}$ forme un système complet d'événements (SCE)[span_4](end_span). On applique la formule des probabilités totales :
$$P(\bar{L}) = P(M_a \cap \bar{L}) + P(M_e \cap \bar{L}) + P(N \cap \bar{L})$$
[span_5](start_span)$$P(\bar{L}) = P(M_a)P(\bar{L}|M_a) + P(M_e)P(\bar{L}|M_e) + P(N)P(\bar{L}|N) \text{[span_5](end_span)}$$
[span_6](start_span)\textbf{Réponse :} $P(\bar{L}) = 0.15 \times 0.5 + 0.08 \times 0.7 + 0.77 \times 0.3 = 0.075 + 0.056 + 0.231 = 0.362$[span_6](end_span).

\subsection*{2. Probabilité que l'auteur soit un prof de maths sachant que le mot est illisible}
[span_7](start_span)\textbf{Raisonnement :} Il s'agit de calculer la probabilité conditionnelle $P(M_a|\bar{L})$ à l'aide de la formule de Bayes[span_7](end_span).
[span_8](start_span)\textbf{Réponse :} $P(M_a|\bar{L}) = \frac{P(M_a \cap \bar{L})}{P(\bar{L})} = \frac{0.075}{0.362} = \frac{75}{362} \approx 0.207$[span_8](end_span).

\hrulefill

\section*{Exercice 2}

\subsection*{1. Démonstration par récurrence}
[span_9](start_span)\textbf{Raisonnement :} On procède par récurrence sur le nombre $n$ d'événements $A_1, A_2, \dots, A_n$[span_9](end_span).
\begin{itemize}
    [span_10](start_span)\item \textbf{Initialisation :} Pour $n=2$, la définition de la probabilité conditionnelle donne bien $P(A_1 \cap A_2) = P(A_1)P(A_2|A_1)$[span_10](end_span).
    [span_11](start_span)\item \textbf{Hérédité :} Supposons la propriété vraie pour un entier $m \ge 2$[span_11](end_span). Pour $m+1$ événements, on écrit :
    [span_12](start_span)$P(A_1 \cap \dots \cap A_{m+1}) = P((A_1 \cap \dots \cap A_m) \cap A_{m+1}) = P(A_1 \cap \dots \cap A_m)P(A_{m+1} | A_1 \cap \dots \cap A_m)$[span_12](end_span).
    En utilisant l'hypothèse de récurrence pour $P(A_1 \cap \dots \cap A_m)$, on obtient directement la formule souhaitée au rang $m+1$. [span_13](start_span)L'hérédité est prouvée[span_13](end_span).
\end{itemize}

\subsection*{2. Tirage dans l'urne}
[span_14](start_span)\textbf{Raisonnement :} Notons $A_k$ l'événement "au $k$-ième tirage, on obtient une boule blanche et une noire" (couleurs différentes) pour $k \in \{1, \dots, n\}$[span_14](end_span).
[span_15](start_span)On cherche $P(A_1 \cap A_2 \cap \dots \cap A_n)$[span_15](end_span).
Au premier tirage, on choisit 2 boules parmi $2n$. Il y a $n^2$ façons de choisir une blanche et une noire :
[span_16](start_span)$P(A_1) = \frac{n^2}{\binom{2n}{2}} = \frac{n^2}{n(2n-1)} = \frac{n}{2n-1}$[span_16](end_span).
Sachant $A_1$, il reste $2n-2$ boules dont $n-1$ de chaque couleur. Donc :
[span_17](start_span)$P(A_2|A_1) = \frac{(n-1)^2}{\binom{2n-2}{2}} = \frac{n-1}{2n-3}$[span_17](end_span).
[span_18](start_span)De manière générale, $P(A_k | A_1 \cap \dots \cap A_{k-1}) = \frac{n-k+1}{2(n-k+1)-1}$[span_18](end_span).
\textbf{Réponse :} En appliquant la formule du 1., on a le produit :
[span_19](start_span)$$P(A_1 \cap \dots \cap A_n) = \frac{n}{2n-1} \times \frac{n-1}{2n-3} \times \dots \times \frac{1}{1} \text{[span_19](end_span)}$$
[span_20](start_span)En multipliant le numérateur et le dénominateur par $(n!) \times 2^n$ pour reconstituer les paires manquantes au dénominateur, on trouve $P = \frac{2^n(n!)^2}{(2n)!}$[span_20](end_span).

\hrulefill

\section*{Exercice 3}
[span_21](start_span)On note $A$ : "l'auteur est anglais" et $V$ : "la lettre tirée est une voyelle"[span_21](end_span).
[span_22](start_span)[span_23](start_span)On a $P(A) = 0.40$ et $P(\bar{A}) = 0.60$[span_22](end_span)[span_23](end_span).
[span_24](start_span)Pour le mot "rigour" (anglais, 6 lettres, 3 voyelles), $P(V|A) = \frac{3}{6} = \frac{1}{2}$[span_24](end_span).
[span_25](start_span)Pour le mot "rigor" (américain, 5 lettres, 2 voyelles), $P(V|\bar{A}) = \frac{2}{5} = 0.4$[span_25](end_span).

\textbf{Raisonnement :} Le système $\{A, \bar{A}\}$ est un SCE. Par la formule des probabilités totales :
[span_26](start_span)$P(V) = P(A)P(V|A) + P(\bar{A})P(V|\bar{A}) = 0.4 \times 0.5 + 0.6 \times 0.4 = 0.44$[span_26](end_span).
[span_27](start_span)On cherche la probabilité $P(A|V)$ via la formule de Bayes : $P(A|V) = \frac{P(A \cap V)}{P(V)}$[span_27](end_span).
[span_28](start_span)\textbf{Réponse :} $P(A|V) = \frac{0.2}{0.44} = \frac{5}{11}$[span_28](end_span).

\hrulefill

\section*{Exercice 4}

\subsection*{1. Famille de trois enfants}
\textbf{Raisonnement :} L'univers est $\Omega = \{F, G\}^3$ avec équiprobabilité, donc $|\Omega| [span_29](start_span)= 8$[span_29](end_span).
$A$ : "enfants des deux sexes" correspond à $\Omega$ privé de $\{(G,G,G), (F,F,F)\}$. Donc $|A| [span_30](start_span)= 6$ et $P(A) = \frac{6}{8} = \frac{3}{4}$[span_30](end_span).
$B$ : "au plus un garçon" correspond à $\{(F,F,F), (G,F,F), (F,G,F), (F,F,G)\}$. Donc $|B| [span_31](start_span)= 4$ et $P(B) = \frac{4}{8} = \frac{1}{2}$[span_31](end_span).
L'intersection $A \cap B$ correspond aux cas "exactement 1 garçon", soit $\{(G,F,F), (F,G,F), (F,F,G)\}$. $|A \cap B| [span_32](start_span)= 3$, donc $P(A \cap B) = \frac{3}{8}$[span_32](end_span).
[span_33](start_span)\textbf{Réponse :} Puisque $P(A) \times P(B) = \frac{3}{4} \times \frac{1}{2} = \frac{3}{8} = P(A \cap B)$, les événements sont \textbf{indépendants}[span_33](end_span).

\subsection*{2. Famille de deux enfants}
\textbf{Raisonnement :} Ici $\Omega = \{F, G\}^2$ avec $|\Omega| [span_34](start_span)= 4$[span_34](end_span).
[span_35](start_span)$A$ = $\{(G,F), (F,G)\}$ d'où $P(A) = \frac{2}{4} = \frac{1}{2}$[span_35](end_span).
[span_36](start_span)$B$ ("au plus un garçon") = $\{(F,F), (G,F), (F,G)\}$ d'où $P(B) = \frac{3}{4}$[span_36](end_span).
[span_37](start_span)L'intersection $A \cap B = \{(G,F), (F,G)\}$, d'où $P(A \cap B) = \frac{1}{2}$[span_37](end_span).
[span_38](start_span)\textbf{Réponse :} $P(A) \times P(B) = \frac{1}{2} \times \frac{3}{4} = \frac{3}{8} \neq P(A \cap B)$, les événements ne sont \textbf{pas indépendants}[span_38](end_span).

\hrulefill

\section*{Exercice 5}
[span_39](start_span)[span_40](start_span)\textbf{Raisonnement :} C'est la répétition de $n$ épreuves de Bernoulli indépendantes (succès = marquer, avec probabilité $p$)[span_39](end_span)[span_40](end_span). [span_41](start_span)L'univers est $\Omega = \{S, E\}^n$[span_41](end_span).
Le nombre de tirages réussis suit une loi binomiale. [span_42](start_span)L'événement "réussir $k$ tirs" correspond aux issues avec $k$ succès et $n-k$ échecs[span_42](end_span). [span_43](start_span)Chaque issue a une probabilité $p^k(1-p)^{n-k}$[span_43](end_span). [span_44](start_span)Il y a $\binom{n}{k}$ façons de placer ces succès[span_44](end_span).
[span_45](start_span)\textbf{Réponse :} La probabilité est $\binom{n}{k}p^k(1-p)^{n-k}$[span_45](end_span).

\hrulefill

\section*{Exercice 6}

\subsection*{1. Probabilité de découvrir le truquage}
\textbf{Raisonnement :} Le jeu possède 33 cartes. L'univers $\Omega$ est l'ensemble des parties à $n$ éléments de ces 33 cartes, $|\Omega| [span_46](start_span)= \binom{33}{n}$[span_46](end_span).
[span_47](start_span)Pour s'apercevoir que le jeu est truqué, il faut tirer \textit{simultanément} les 2 dames de cœur[span_47](end_span). [span_48](start_span)Le nombre de tirages contenant ces 2 cartes est le choix des $n-2$ autres cartes parmi les 31 restantes, soit $\binom{31}{n-2}$[span_48](end_span).
[span_49](start_span)\textbf{Réponse :} $P(truqué) = \frac{\binom{31}{n-2}}{\binom{33}{n}} = \frac{31!}{(n-2)!(33-n)!} \times \frac{n!(33-n)!}{33!} = \frac{n(n-1)}{33 \times 32}$[span_49](end_span).

\subsection*{2. Nombre d'expériences avec $n=4$}
[span_50](start_span)\textbf{Raisonnement :} Pour $n=4$, $P(truqué) = \frac{4 \times 3}{33 \times 32} = \frac{1}{88}$[span_50](end_span).
[span_51](start_span)En répétant $k$ fois, la probabilité de ne \textit{jamais} s'en apercevoir est $(1 - \frac{1}{88})^k = (\frac{87}{88})^k$[span_51](end_span).
[span_52](start_span)On cherche $k$ tel que $1 - (\frac{87}{88})^k \ge 0.95$, soit $(\frac{87}{88})^k \le 0.05$[span_52](end_span).
[span_53](start_span)\textbf{Réponse :} En passant au logarithme : $k \ln(87/88) \le \ln(0.05) \implies k \ge \frac{\ln(0.05)}{\ln(87/88)}$ (le sens de l'inégalité change car $\ln(87/88) < 0$)[span_53](end_span). [span_54](start_span)On trouve $k \ge 263$[span_54](end_span).

\hrulefill

\section*{Exercice 7}
[span_55](start_span)On note $V$ : "vacciné" et $M$ : "tombe malade"[span_55](end_span).
[span_56](start_span)[span_57](start_span)L'énoncé donne $P(V) = \frac{1}{4}$ (donc $P(\bar{V}) = \frac{3}{4}$) et $P(M|V) = \frac{1}{20}$[span_56](end_span)[span_57](end_span). 
[span_58](start_span)[span_59](start_span)De plus, parmi les malades, 4 sont non-vaccinés pour 1 vacciné, ce qui signifie que $P(\bar{V}|M) = \frac{4}{5}$ et $P(V|M) = \frac{1}{5}$[span_58](end_span)[span_59](end_span).
[span_60](start_span)\textbf{Raisonnement :} On cherche $P(M|\bar{V})$[span_60](end_span).
[span_61](start_span)On sait que $P(V|M) = \frac{P(M \cap V)}{P(M)} = \frac{P(M|V)P(V)}{P(M \cap V) + P(M \cap \bar{V})} = \frac{P(M|V)P(V)}{P(M|V)P(V) + P(M|\bar{V})P(\bar{V})}$[span_61](end_span).
[span_62](start_span)En injectant les valeurs : $\frac{1}{5} = \frac{\frac{1}{20} \times \frac{1}{4}}{\frac{1}{20} \times \frac{1}{4} + P(M|\bar{V}) \times \frac{3}{4}} = \frac{\frac{1}{80}}{\frac{1}{80} + \frac{3}{4}P(M|\bar{V})}$[span_62](end_span).
[span_63](start_span)On obtient : $\frac{1}{80} + \frac{3}{4}P(M|\bar{V}) = \frac{5}{80} \implies \frac{3}{4}P(M|\bar{V}) = \frac{4}{80} = \frac{1}{20}$[span_63](end_span).
[span_64](start_span)\textbf{Réponse :} $P(M|\bar{V}) = \frac{1}{20} \times \frac{4}{3} = \frac{1}{15}$[span_64](end_span).

\hrulefill

\section*{Exercice 8}
[span_65](start_span)[span_66](start_span)On note $M$ : "être malade" ($P(M) = 0.01$), $T$ : "le test est positif", et $D$ : "décéder"[span_65](end_span)[span_66](end_span).
[span_67](start_span)On a $P(T|M) = 0.8$, $P(\bar{T}|M) = 0.2$, $P(T|\bar{M}) = 0.03$, $P(\bar{T}|\bar{M}) = 0.97$[span_67](end_span).
Probabilités de décès selon le statut et le traitement (on est traité ssi le test est positif) :
[span_68](start_span)$P(D|M \cap \bar{T}) = 0.5$ (malade non traité)[span_68](end_span)[span_69](start_span)[span_70](start_span), $P(D|M \cap T) = 0.1$ (malade traité)[span_69](end_span)[span_70](end_span),
[span_71](start_span)[span_72](start_span)$P(D|\bar{M} \cap T) = 0.02$ (sain traité)[span_71](end_span)[span_72](end_span)[span_73](start_span), $P(D|\bar{M} \cap \bar{T}) = 0$ (sain non traité)[span_73](end_span).

\subsection*{1. Sans test de dépistage}
[span_74](start_span)[span_75](start_span)\textbf{Raisonnement :} Le décès ne survient que chez les malades avec probabilité $0.5$[span_74](end_span)[span_75](end_span).
[span_76](start_span)\textbf{Réponse :} $P(D) = P(M \cap D) = P(M)P(D|M) = 0.01 \times 0.5 = 0.005$[span_76](end_span).

\subsection*{2. Avec dépistage généralisé}
[span_77](start_span)\textbf{Raisonnement :} La famille $\{M \cap T, M \cap \bar{T}, \bar{M} \cap T, \bar{M} \cap \bar{T}\}$ forme un SCE[span_77](end_span). 
[span_78](start_span)La formule des probabilités totales donne[span_78](end_span):
[span_79](start_span)$P(D) = P(M)P(T|M)P(D|M \cap T) + P(M)P(\bar{T}|M)P(D|M \cap \bar{T}) + P(\bar{M})P(T|\bar{M})P(D|\bar{M} \cap T) + 0$[span_79](end_span)
[span_80](start_span)\textbf{Réponse :} $P(D) = 0.01 \times 0.8 \times 0.1 + 0.01 \times 0.2 \times 0.5 + 0.99 \times 0.03 \times 0.02 = 0.0008 + 0.001 + 0.000594 = 0.002394$[span_80](end_span).

\hrulefill

\section*{Exercice 9}
\textbf{Raisonnement :} Il y a 9 personnes au total. On choisit 4 personnes. [span_81](start_span)Le nombre de cas possibles est $\binom{9}{4} = 126$[span_81](end_span).
[span_82](start_span)L'événement "choisir 4 femmes" n'est possible que d'une seule façon, car il n'y a que 4 femmes dans le groupe (on les prend toutes) : $\binom{4}{4} = 1$[span_82](end_span).
[span_83](start_span)\textbf{Réponse :} $P = \frac{1}{126}$[span_83](end_span).

\hrulefill

\section*{Exercice 10}
\textbf{Raisonnement :} Notons $S_1, S_2, S_3$ les événements "l'alliance est dans le secteur $i$". [span_84](start_span)[span_85](start_span)À l'origine, ces probabilités sont équiprobables : $P(S_1) = P(S_2) = P(S_3) = \frac{1}{3}$[span_84](end_span)[span_85](end_span).
[span_86](start_span)Notons $E$ l'événement "les fouilles sur le secteur 1 ne donnent rien"[span_86](end_span).
[span_87](start_span)D'après l'énoncé, $P(E|S_1) = p$ (probabilité de rater l'alliance si elle y est)[span_87](end_span). [span_88](start_span)Si l'alliance est ailleurs, on est sûr de ne rien trouver dans le secteur 1 : $P(E|S_2) = P(E|S_3) = 1$[span_88](end_span).
On calcule la probabilité de $E$ par les probabilités totales :
[span_89](start_span)$P(E) = P(E|S_1)P(S_1) + P(E|S_2)P(S_2) + P(E|S_3)P(S_3) = p(\frac{1}{3}) + 1(\frac{1}{3}) + 1(\frac{1}{3}) = \frac{p+2}{3}$[span_89](end_span).
On utilise ensuite le théorème de Bayes pour réévaluer les probabilités :
[span_90](start_span)\textbf{Réponses :}[span_90](end_span)
\begin{itemize}
    [span_91](start_span)\item $P(S_1|E) = \frac{P(E|S_1)P(S_1)}{P(E)} = \frac{p/3}{(p+2)/3} = \frac{p}{p+2}$[span_91](end_span)
    [span_92](start_span)\item $P(S_2|E) = \frac{P(E|S_2)P(S_2)}{P(E)} = \frac{1/3}{(p+2)/3} = \frac{1}{p+2}$[span_92](end_span)
    [span_93](start_span)\item $P(S_3|E) = \frac{P(E|S_3)P(S_3)}{P(E)} = \frac{1/3}{(p+2)/3} = \frac{1}{p+2}$[span_93](end_span)
\end{itemize}

\hrulefill

\section*{Exercice 11}
\subsection*{1. Expressions des probabilités au jour $n+1$}
\textbf{(a) Raisonnement :} Par la formule des probabilités totales avec le SCE $\{M_n, S_n, B_n\}$ :
[span_94](start_span)[span_95](start_span)$P(M_{n+1}) = P(M_{n+1}|M_n)P(M_n) + P(M_{n+1}|S_n)P(S_n) + P(M_{n+1}|B_n)P(B_n)$[span_94](end_span)[span_95](end_span).
[span_96](start_span)D'après les règles de transition données[span_96](end_span):
\textbf{Réponses :}
[span_97](start_span)[span_98](start_span)$p_{n+1} = (1-2a)p_n + a q_n + a r_n$[span_97](end_span)[span_98](end_span).
[span_99](start_span)[span_100](start_span)$q_{n+1} = a p_n + (1-2a)q_n + a r_n$[span_99](end_span)[span_100](end_span).

\textbf{(b) Raisonnement :} À tout jour $n$, le titre doit soit monter, soit être stable, soit baisser. [span_101](start_span)Donc la somme des probabilités vaut 1[span_101](end_span).
\textbf{Réponse :} $p_n + q_n + r_n = 1$. [span_102](start_span)On en déduit $r_n = 1 - p_n - q_n$[span_102](end_span).

\subsection*{2. Suites arithmético-géométriques}
\textbf{Raisonnement :} On injecte l'expression de $r_n$ dans $p_{n+1}$ et $q_{n+1}$ :
[span_103](start_span)$p_{n+1} = (1-2a)p_n + a q_n + a(1 - p_n - q_n) = (1-2a)p_n - a p_n + a = (1-3a)p_n + a$[span_103](end_span).
[span_104](start_span)De même, $q_{n+1} = a p_n + (1-2a)q_n + a(1 - p_n - q_n) = (1-3a)q_n + a$[span_104](end_span).
[span_105](start_span)\textbf{Réponse :} Les suites $(p_n)$ et $(q_n)$ s'expriment sous la forme $u_{n+1} = \alpha u_n + \beta$, elles sont donc arithmético-géométriques de raison $1-3a$[span_105](end_span).

\subsection*{3. Expression explicite et limites}
[span_106](start_span)\textbf{Raisonnement :} Le point fixe de ces suites est la solution de $\ell = (1-3a)\ell + a \implies 3a\ell = a \implies \ell = \frac{1}{3}$ (puisque $a \neq 0$)[span_106](end_span).
On pose la suite auxiliaire $u_n = p_n - 1/3$ (resp. $v_n = q_n - 1/3$) qui est géométrique de raison $1-3a$. 
Donc $p_n = \frac{1}{3} + (p_1 - \frac{1}{3})(1-3a)^{n-1}$.
[span_107](start_span)L'énoncé stipule qu'au premier jour le titre est stable, donc $p_1 = 0$, $q_1 = 1$, $r_1 = 0$[span_107](end_span).
\textbf{Réponses :}
[span_108](start_span)$p_n = \frac{1}{3} - \frac{1}{3}(1-3a)^{n-1}$[span_108](end_span).
[span_109](start_span)$q_n = \frac{1}{3} + \frac{2}{3}(1-3a)^{n-1}$[span_109](end_span).
[span_110](start_span)$r_n = 1 - p_n - q_n = \frac{1}{3} - \frac{1}{3}(1-3a)^{n-1}$[span_110](end_span).

[span_111](start_span)\textbf{Limites :} Puisque $a \in ]0, 1/2[$, on a $1-3a \in ]-1/2, 1[$[span_111](end_span). [span_112](start_span)La suite $(1-3a)^{n-1}$ tend donc vers 0 quand $n \to +\infty$[span_112](end_span).
[span_113](start_span)Les limites des trois suites sont $\lim_{n \to \infty} p_n = \frac{1}{3}$, $\lim_{n \to \infty} q_n = \frac{1}{3}$, $\lim_{n \to \infty} r_n = \frac{1}{3}$[span_113](end_span).
[span_114](start_span)\textbf{Interprétation :} À long terme, le système "oublie" sa condition initiale et chaque comportement (monter, stagner, baisser) devient équiprobable[span_114](end_span).

\end{document}
